\documentclass[11pt,a4paper]{article}
\usepackage[margin=25mm]{geometry}
\usepackage[T1]{fontenc}
\usepackage{lmodern}
\usepackage{amsmath,amssymb,booktabs,array}
\usepackage[hidelinks]{hyperref}
\usepackage{microtype}
\setlength{\parindent}{0pt}
\setlength{\parskip}{6pt}
\emergencystretch=2em
\newcommand{\BR}{\operatorname{BR}}
\newcommand{\E}{\mathbb{E}}
\newcommand{\Prb}{\mathbb{P}}
\newcommand{\sourcepage}[1]{\textit{Original notes: #1.}}
\title{An Exploratory Model of Anticipation\\
\large Edited model-background edition with clarifications}
\author{Researcher's original model-building notes\\
\normalsize Edited with AI assistance for private mentor review}
\date{24 October 2026}
\begin{document}
\maketitle

\textbf{Status and purpose.} This edited model-background edition presents the
researcher's mathematical approach from early iPad notes, with explicit
clarifications. It is not a verbatim transcription or a final empirical protocol.
The full eight-page-source reading is preserved in the local v3 backup; this
review edition deliberately abbreviates the retired measurement proposal.

Investors may regard familiar daily moving averages as support and anticipate
others acting before contact. The model expresses that intuition as choices,
beliefs, a payoff trade-off and a best-response fixed point. It is not a claim
of theoretical novelty or an empirically identified mechanism. The principal
empirical deliverable remains the final measured distribution of real
successful-rebound pivots; the current mainline review note is authoritative
for measurement. Source-page references identify where the mathematical
arguments originated, not current implementation instructions.

\section{Distance, attention and beliefs}
\sourcepage{page 1.}

Let $P_t$ denote price and $M_t>0$ a technical reference. Define
\begin{equation}
 x_t=\frac{P_t-M_t}{M_t}.
\end{equation}
Thus $x_t=0.003$ means $0.3\%$ above the reference, $x_t=0$ is on it,
and $10{,}000x_t$ is signed distance in basis points. The draft posits
\begin{equation}
 A(x)>0,\qquad A'(x)<0
\end{equation}
on the relevant approach domain. The interpretation is greater attention nearer
the line from above. This is an assumption, not an attention measure observed
in the price data.

Trader $i$ chooses advance-entry distance $d_i\geq0$. For $N$ traders, write
\begin{equation}
 d_{-i}=(d_1,\ldots,d_{i-1},d_{i+1},\ldots,d_N),\quad
 \bar d_{-i}=\frac{1}{N-1}\sum_{j\ne i}d_j,\quad
 \mu_i=\E_i[\bar d_{-i}].
\end{equation}
Here $\E_i$ reflects trader $i$'s information and beliefs. For example, with four
traders, $\bar d_{-2}=(d_1+d_3+d_4)/3$. Summarizing the distribution of other
traders' actions and beliefs by one mean is a substantive simplification.

The proposed best-response relation and complementarity assumption are
\begin{equation}
 d_i=\BR(\mu_i;A),\qquad \BR_\mu>0.
\end{equation}
A symmetric equilibrium with identical relevant primitives and consistent beliefs
must satisfy
\begin{equation}\label{eq:fixedpoint}
 d^*=\BR(d^*;A).
\end{equation}
A scalar fixed point is not by itself a complete heterogeneous market model.
Complementarity requires its own assumptions; it will not follow merely from
attention increasing near the line.

\section{Entry gain, capture probability and cost}
\sourcepage{pages 2--4.}

For a single stylized decision, \emph{hold the reference $M>0$ and rebound sale
price $p^R$ fixed}. Choosing $d$ then gives
\begin{equation}
 P^{\mathrm{entry}}=M(1+d),\qquad
 G(d)=p^R-M(1+d),\qquad G'(d)=-M,\qquad G''(d)=0.
\end{equation}
Entering earlier at a higher price reduces the gain conditional on capture.
If the model requires $G(d)\geq0$, its feasible domain must impose
$d\leq p^R/M-1$, provided that bound is nonnegative. The draft's proposed
finite upper bound on optimal entry distance needs conditions on the
\emph{whole payoff}, not just $G'<0$.

Let $q(d;\mu,A)\in[0,1]$ be a capture probability and let $C(d)$ represent
early-entry cost or risk. The draft assumes
\begin{equation}
 q_d>0,\qquad C'(d)>0
\end{equation}
and proposes the reduced-form expected payoff
\begin{equation}\label{eq:payoff}
 U(d,\mu;A)=q(d;\mu,A)G(d)-C(d).
\end{equation}
Failed capture has no separate payoff here; such a loss would need to be
modeled explicitly or included in $C$. These formulas are not a validated
trading-profit model. In particular, differentiating with fixed $M$ cannot
silently be applied to a live MA that itself changes with price.

For an \emph{interior}, differentiable maximizer, holding the specified attention
environment fixed in the notation, the first-order condition is
\begin{align}
 U_d&=q_dG+qG'-C'=q_dG-Mq-C'=0, \label{eq:foc}\\
 U_{dd}&=q_{dd}G+2q_dG'+qG''-C''
       =q_{dd}G-2Mq_d-C''. \label{eq:soc}
\end{align}
The first term in \eqref{eq:foc} is the marginal capture benefit; the other
terms reflect the higher entry price and marginal cost. At an interior symmetric
solution, evaluate this condition at $d=\mu=d^*$.

\textbf{Clarification of page 4.} At a stationary point, $U_{dd}<0$ is a
sufficient condition for a strict local maximum, not a necessary condition for
every optimum or equilibrium. Nor does it alone establish a global or unique
best response. With suitable regularity, an interior local maximum satisfies
$U_d=0$ and $U_{dd}\leq0$; a zero second derivative may require higher-order
analysis. Boundary candidates must be checked separately. Strict concavity on
a compact feasible interval, together with continuity, is one possible stronger
route to an attained, unique maximizer.

\section{Existence, local approximation and stability}
\sourcepage{pages 1--5.}

\subsection{A fixed-point existence argument}
Suppose a continuous, single-valued best response is a self-map
\[
 \BR:[0,\bar d]\longrightarrow[0,\bar d],\qquad 0<\bar d<\infty.
\]
For $F(d)=\BR(d)-d$, the endpoint inequalities are
\[
 F(0)=\BR(0)\geq0,\qquad F(\bar d)=\BR(\bar d)-\bar d\leq0.
\]
The intermediate value theorem gives some $d^*\in[0,\bar d]$ with
$F(d^*)=0$. The endpoint signs need not be strict. This establishes existence,
possibly at a boundary, \emph{given} the self-map and continuity assumptions.
It does not establish uniqueness, a positive interior equilibrium or stability.
The notes do not derive all these best-response properties from payoff primitives.

\subsection{What strategic complementarity requires}
Along a differentiable interior best-response branch with $U_{dd}\ne0$,
implicit differentiation of $U_d(\BR(\mu),\mu;A)=0$ gives
\begin{equation}
 \BR_\mu=-\frac{U_{d\mu}}{U_{dd}},\qquad
 U_{d\mu}=Gq_{d\mu}-Mq_\mu. \label{eq:brslope}
\end{equation}
The second expression uses fixed $M$, fixed $p^R$ and $C=C(d)$. If
$U_{dd}<0$, positive complementarity requires $U_{d\mu}>0$. Neither
$q_d>0$ nor $A'(x)<0$ alone implies this cross-effect.

\subsection{Taylor approximation versus an affine specification}
Around a reference belief $\mu_0$, differentiability gives
\begin{align}
 \BR(\mu)&=\BR(\mu_0)+\BR_\mu(\mu_0)(\mu-\mu_0)
              +o(|\mu-\mu_0|)\\
         &\approx \alpha+\beta\mu,\qquad
 \beta=\BR_\mu(\mu_0),\quad
 \alpha=\BR(\mu_0)-\beta\mu_0.
\end{align}
\textbf{Clarification of page 2.} The approximation is local. Writing it as
an exact global equality discards the remainder and introduces an additional
affine-model assumption. If $\BR(\mu)=\alpha+\beta\mu$ really is affine,
then, for $\beta\ne1$,
\begin{equation}
 d^*=\frac{\alpha}{1-\beta},\qquad
 \BR(\mu)-d^*=\beta(\mu-d^*).
\end{equation}
These equalities must also respect the feasible action interval.

\subsection{A hypothetical best-response iteration}
Consider $d_{n+1}=\BR(d_n;A)$ and perturb an interior fixed point:
$d_n=d^*+\delta_n$. With $\beta=\BR_\mu(d^*)$,
\begin{equation}
 \delta_{n+1}=\beta\delta_n+o(|\delta_n|). \label{eq:local}
\end{equation}
If the required second-order smoothness holds, the remainder can instead be
written as $O(\delta_n^2)$. The linearized model gives
\[
 \delta_n^{\mathrm{lin}}=\beta^n\delta_0.
\]
For a genuinely affine best response this is exact for the actual iterates.
For a general nonlinear response it is not an exact formula for all $n$.

Under standard smoothness conditions, $|\beta|<1$ is sufficient for local
convergence; $|\beta|>1$ is locally unstable, while $|\beta|=1$ requires
further analysis. A uniform contraction on the entire closed action interval
is a stronger sufficient condition for a unique fixed point and convergence
from every starting point in that interval.

\textbf{Correction of page 1 and qualification of page 5.} A unique equilibrium
does \emph{not} imply $|\BR_\mu|<1$. Existence, uniqueness and stability
are distinct. The power formula on page 5 is exact only for the affine or
linearized model. Moreover, the iteration index $n$ is not automatically real
market time; these dynamics were not estimated from the price data.

\section{Positive anticipation and salience}
\sourcepage{pages 5--6.}

\subsection{A condition ruling out zero as an optimum}
At $d=\mu=0$, the draft considers
\begin{equation}
 U_d(0,0;A)=q_d(0,0;A)G(0)-Mq(0,0;A)-C'(0)>0.
\end{equation}
With feasible points immediately to the right of zero and the relevant right
derivative, this rules out zero as a maximizing choice when $\mu=0$.
If a global maximum is attained, any maximizing choice must then be positive.
With the continuous-self-map existence conditions above and $\BR(0)>0$,
a fixed point cannot equal zero; its value might still be the upper boundary.
This is a conditional argument, not an unconditional positive interior result.

For the affine benchmark, $\alpha>0$ and $0\leq\beta<1$ yield a positive
$d^*=\alpha/(1-\beta)$ only if it lies in the action domain and the assumed
self-map conditions are met. A positive local Taylor intercept alone does not
establish that conclusion for an arbitrary nonlinear response.

\subsection{A conditional comparative-static formula}
Let salience enter through $A(x;s)$ and write the resulting payoff as
$U(d,\mu;s)$. Along a differentiable interior equilibrium branch,
\[
 U_d(d^*(s),d^*(s);s)=0.
\]
Differentiation gives
\begin{align}
 (U_{dd}+U_{d\mu})\frac{d d^*}{ds}+U_{ds}&=0,\\
 \frac{d d^*}{ds}
 &= -\frac{U_{ds}}{U_{dd}+U_{d\mu}}
  = -\frac{U_{ds}}{U_{dd}(1-\beta)},\qquad
 \beta=-\frac{U_{d\mu}}{U_{dd}}. \label{eq:salience}
\end{align}
All derivatives are evaluated at the equilibrium and include the specified
chain-rule effects of $s$. The branch must be nondegenerate:
$U_{dd}+U_{d\mu}\ne0$. The last expression cannot be used at $\beta=1$.
Under $U_{dd}<0$ and $\beta<1$, the direction is the sign of $U_{ds}$,
which needs an additional substantive assumption. Greater salience does not
unconditionally imply earlier action. No payoff derivatives in
\eqref{eq:salience} were structurally estimated by the completed study.

\section{Reaction distance and a proposed CDF bridge}
\sourcepage{pages 6--7.}

The draft introduces $D_R$ as the random distance from the reference at which
a reaction begins. In an idealized approach from above, entry at threshold $d$
before reversal is proposed to correspond to $D_R\leq d$. It therefore writes
\begin{equation}
 q(d;\mu,A)=H(d\mid\mu,A)
          =\Prb(D_R\leq d\mid\mu,A).
\end{equation}
This needs an event-equivalence assumption: the path reaches the entry threshold
before reversing precisely under the stated distance ordering, with compatible
reference and execution assumptions. Gaps, moving references, repeated
approaches, non-reactions and finite windows complicate the equivalence.
How non-reactions enter the model and which population is conditioned on must
be specified.

If $H$ is appropriately smooth, let $h=\partial H/\partial d$. The payoff
and interior first-order condition can then be written
\begin{align}
 U(d,\mu;A)&=H(d\mid\mu,A)G(d)-C(d),\\
 h(d\mid\mu,A)G(d)+H(d\mid\mu,A)G'(d)-C'(d)&=0.
\end{align}
At a symmetric interior equilibrium, set $d=\mu=d^*$.
A CDF is nondecreasing; a strictly positive density is an additional
assumption, not a generic CDF property.

For observed reaction distances $D_{R,1},\ldots,D_{R,m}$, the notes propose
\begin{equation}
 \widehat H(d)=\frac1m\sum_{j=1}^{m}\mathbf{1}\{D_{R,j}\leq d\}.
\end{equation}
An empirical CDF is a step function, not an ordinary smooth density.
More fundamentally, if observations are selected because they are successful
rebounds, this estimates a \emph{successful-event location distribution}.
It does not identify the capture probability across all opportunities or
the probability of rebound conditional on distance.

The theoretical trader choice $d_i$, equilibrium $d^*$ and empirically measured
pivot distance are distinct objects. A median, mode or KDE peak of the
successful-event distribution cannot be substituted for $d^*$.

\section{The early measurement proposal and its later replacement}
\label{sec:oldmeasurement}
\sourcepage{pages 7--8.}

\subsection{Live daily reference}
Page 7 proposes
\begin{equation}
 M_t^{(5)}=\frac{P_t+C_{t-1}+C_{t-2}+C_{t-3}+C_{t-4}}5,
\end{equation}
where $C_{t-k}$ denotes the close of the $k$th previous completed trading day,
not the previous minute. The motivation is to approximate a currently
observable technical reference. It is not evidence that every investor uses
that display. A reference built only from the five previous completed daily
closes is a different benchmark, proposed in the notes as a possible robustness
comparison rather than substituted into this live formula.

\subsection{Current measurement takes precedence}
Page 8 proposed an earlier band-entry origin and running-low threshold detector.
Its detailed equations are omitted from this edited background edition because
they are not the final empirical rules. The complete historical reading remains
in the local v3 backup. The final sample instead uses the frozen bullish-pullback
A background, original-near origins, full EWMA, three consecutive qualifying
closes within five minutes, and the fixed non-overlap rule. Its pivot is the last
tied lowest close before the first qualifying three-close window. See
\texttt{MAINLINE\_REVIEW.md} for the current definitions and empirical results.
These are conditional observations, not estimates of the model's optimal
choice or equilibrium distance.

\section*{Source and editorial scope}
This edited reading retains the mathematical model and clarifications from \emph{Project notes.pdf}, supplied by the researcher. The original
source's SHA-256 is recorded in the candidate's provenance manifest.
Corrections are identified in the text, including fixed-reference derivatives,
interior versus boundary optima, fixed-point existence versus uniqueness and
stability, local versus exact affine formulas, conditional comparative statics,
and the empirical-selection gap in the CDF interpretation.

No original research report, data, analysis code or frozen protocol has been
changed. No new empirical experiment or theory-validation claim is introduced.
The handwritten source and complete historical transcription remain local; this edited reading and its source are
the intended review appendix. Research and GitHub publication remain stopped
pending the owner's separate instructions.
\end{document}